\documentclass[10pt,twocolumn]{article}
% =====================================================================
%  preamble.tex — shared preamble for the agent-endpoint defense paper
%  Compile with pdflatex/latexmk or tectonic:  tectonic main.tex
% =====================================================================

% ---------- Page geometry & fonts ----------
\usepackage[a4paper,margin=0.75in,columnsep=0.28in]{geometry}
\usepackage[T1]{fontenc}
\usepackage[utf8]{inputenc}
\usepackage{newtxtext,newtxmath}      % Times-like text + matching math
\usepackage[scaled=0.85]{beramono}    % readable monospace for code/paths
\usepackage{microtype}                % better justification/protrusion

% ---------- Math, units, numbers ----------
\usepackage{amsmath}
\usepackage{siunitx}
\sisetup{group-separator={,},group-minimum-digits=4,detect-all}

% ---------- Tables ----------
\usepackage{booktabs}
\usepackage{multirow}
\usepackage{tabularx}
\usepackage{makecell}
\renewcommand{\arraystretch}{1.12}

% ---------- Figures & plots ----------
\usepackage{graphicx}
\usepackage{xcolor}
\usepackage{tikz}
\usetikzlibrary{arrows.meta,positioning,fit,backgrounds,shapes.geometric}
\usepackage{pgfplots}
\pgfplotsset{compat=1.18}
\usepackage[font=small,labelfont=bf,skip=4pt]{caption}
\usepackage{subcaption}

% ---------- Lists, code, misc ----------
\usepackage{enumitem}
\setlist{nosep,leftmargin=*}
\usepackage{listings}
\lstset{basicstyle=\ttfamily\footnotesize,columns=fullflexible,
        breaklines=true,frame=single,framerule=0.3pt,rulecolor=\color{black!30},
        xleftmargin=2pt,xrightmargin=2pt,aboveskip=4pt,belowskip=4pt}
\usepackage{balance}                  % balance last-page columns

% ---------- Colour palette ----------
\definecolor{accent}{HTML}{1F4E79}
\definecolor{good}{HTML}{2E7D32}
\definecolor{bad}{HTML}{B71C1C}
\definecolor{warn}{HTML}{E65100}
\definecolor{muted}{HTML}{5F6368}

% ---------- Section styling ----------
\usepackage{titlesec}
\titleformat{\section}{\normalfont\large\bfseries\color{accent}}{\thesection}{0.6em}{}
\titleformat{\subsection}{\normalfont\normalsize\bfseries}{\thesubsection}{0.5em}{}
\titlespacing*{\section}{0pt}{1.1ex plus .3ex}{0.6ex}
\titlespacing*{\subsection}{0pt}{0.9ex plus .2ex}{0.4ex}

% ---------- Bibliography ----------
\usepackage[numbers,sort&compress]{natbib}
\setlength{\bibsep}{2pt plus 0.3ex}

% ---------- Hyperlinks (load late) ----------
\usepackage[hidelinks,pdfusetitle]{hyperref}
\usepackage{url}
\usepackage[capitalise,noabbrev]{cleveref}

% ---------- Project macros ----------
\newcommand{\sys}{\textsc{Guardian}}
\newcommand{\kit}{\texttt{agentkit}}
\newcommand{\code}[1]{\texttt{\small #1}}
\newcommand{\us}{\,\si{\micro\second}}
\newcommand{\ms}{\,\si{\milli\second}}
\newcommand{\pass}{\textcolor{good}{\ding{51}}}
\newcommand{\fail}{\textcolor{bad}{\ding{55}}}
\usepackage{pifont}
% Red, unmissable placeholder — search for "\placeholder" before submitting.
\newcommand{\placeholder}[1]{\textcolor{bad}{\textbf{[#1]}}}
\newcommand{\finding}[1]{\par\smallskip\noindent\textbf{\textcolor{accent}{#1}}\ }


% Real-traffic numbers are generated from the exposure study (artifact/gen_realtraffic_tex.py).
\newif\ifrealdata
\IfFileExists{realtraffic.tex}{\input{realtraffic}\realdatatrue}{\realdatafalse}

\title{\Large\bfseries Guarding the Door: Measuring and Hardening Lightweight\\
Application-Layer Defense for an Internet-Exposed {MCP} Agent Endpoint}
\author{%
  Niranjan R.\\
  Department of Artificial Intelligence and Data Science\\
  KPR Institute of Engineering and Technology (KPRIET), Arasur, Coimbatore\\
  \texttt{niranjan.r.dev@proton.me}
}
\date{}

\begin{document}
\maketitle

% =====================================================================
\begin{abstract}
Tool-using AI agents are increasingly reached over the network through the
Model Context Protocol (MCP), which turns a personal automation stack into an
Internet-facing HTTP service. We study what a single operator can realistically
put in front of such an endpoint. We evaluate \sys{}, a standard-library-only
admission gate deployed in front of a self-hosted 29-tool MCP gateway published
through a Cloudflare tunnel. The evaluation uses a reproducible harness of 16
attack and benign scenarios with latency and memory benchmarks, run across three
revisions: the unmodified original (v1), a hardened revision (v2), and a
post-authentication extension (v3). v1 had no false
blocks on benign MCP sessions and stopped declared crawlers, single-source
credential probing and floods at \SI{5}{\micro\second} median overhead. It also
had measurable weaknesses, all caused by keying state on individual client IPs:
distributed guessing went unbanned (700 guesses, 0 bans), ban-list lookups grew
linearly (\SI{21.7}{\milli\second} per request at 50k bans), rate-limit memory
was unbounded (\SI{87}{MB} at 100k sources), and there were further gaps in IPv6
rotation, header trust and lockout. v2 closes each of these with a small local
change. Latency stays flat at \SI{0.13}{\milli\second} up to 50k bans
($169\times$ faster), memory is capped (\SI{18.8}{MB} at 100k sources),
distributed guessing earns 95 bans, IPv6 rotation inside a /64 costs one ban
instead of 200, and a mistyped token costs 5 minutes rather than an hour. None of
this introduces benign false blocks, and v2 adds \SI{9}{\micro\second} per
request. v3 extends defense past the admission layer: scoped bearer tokens
limit each credential to a named set of tools, and per-tool rate limiting
throttles authorized abuse of individual tools without affecting others. Four
concurrent clients complete full MCP sessions with no cross-contamination,
and long-running sessions recover cleanly across rate windows. User-agent
spoofing remains undetectable at this layer; the bearer token still stops it.
\ifrealdata
Over a \RTHours-hour public exposure (\RTRequests{} requests from \RTSources{}
sources), replaying the same real trace through each policy shows v1 blocking
\RTvOneBlocked{} requests at the gate and v2 blocking \RTvTwoBlocked{}.
\else
A time-boxed public exposure study is in progress; its results will be
replayed through all three policies using the validated trace-replay tooling.
\fi
We release the harness, trace-replay tooling and aggregate results.
\end{abstract}

% =====================================================================
\section{Introduction}
Agent frameworks let a language model call tools~\citep{yao2023react,schick2023toolformer,patil2023gorilla},
and MCP standardizes how those tools are exposed over stdio or HTTP~\citep{mcpspec}.
An MCP server published on a public hostname inherits every property of an
Internet service. Internet-wide scanners find it~\citep{durumeric2013zmap}. AI
data collectors crawl it, and their volume and opt-out behaviour are in
flux~\citep{longpre2024consent,rfc9309}. Others probe it for credentials. Unlike
a static website, the thing behind the door can \emph{act}. The security
community has responded: surveys catalogue tool poisoning and over-privileged
servers~\citep{hou2025mcp,owaspllm2025}; benchmarks quantify tool-poisoning
success rates above 72\%~\citep{mcptox2025}; and the NSA has published formal
design guidance~\citep{nsamcp2026}. Yet a recent systematic review of 73
defense studies~\citep{tamayo2026slr} finds that all existing work targets
post-authentication threats. We ask the question that comes first:
\emph{who can reach the endpoint at all, and at what cost to the operator?}
Enterprises answer with identity-aware proxies and managed bot defense.
Individual developers who self-host an MCP server often have only the server
process itself.

\begin{enumerate}
  \item[\textbf{RQ1}] Does a small, dependency-free admission gate stop the
        unwanted traffic an exposed agent endpoint attracts, without blocking
        legitimate MCP clients?
  \item[\textbf{RQ2}] Where, measurably, does per-client-IP defense break down?
  \item[\textbf{RQ3}] Can those failures be fixed locally, and at what cost?
  \item[\textbf{RQ4}] How do the policies compare on real Internet traffic?
\end{enumerate}

\noindent\textbf{Contributions.}
(i) A deployed defense for an MCP endpoint that combines an admission gate,
bearer authentication and capability minimization (\cref{sec:design}).
(ii) A reproducible evaluation of the unmodified code, with eight quantified
weaknesses (\cref{sec:v1}).
(iii) A hardened revision and a before/after comparison on identical scenarios
(\cref{sec:v2}).
(iv) A post-authentication extension with scoped tokens and per-tool rate
limiting, evaluated on two new scenarios (\cref{sec:v3}).
(v) A time-boxed public exposure whose trace is replayed through every
policy (\cref{sec:real}).
(vi) Release of the harness, replay tooling and aggregate data.

\noindent\textbf{Claim.} \emph{For a single-operator, Internet-exposed MCP
endpoint, a lightweight application-layer gate combined with a
high-entropy bearer token and a minimized tool surface stops declared
crawlers, credential probing, path scanning and floods at microsecond
cost with no false blocks on benign MCP sessions. It is effective only
if its state is keyed on network sources rather than individual addresses
and bounded in size. Scoped tokens and per-tool rate limits extend the
defense past authentication, bounding the blast radius of a leaked
credential. The token, not the gate, remains the security boundary.}

% =====================================================================
\section{System and Threat Model}\label{sec:design}

\subsection{Deployment}
The protected service is \kit{}, a personal tool gateway that exposes one
29-tool catalog (DevOps monitoring, a shell assistant, a personal-memory
agent, a breach monitor and a homelab reconnaissance agent) over stdio and
MCP Streamable HTTP~\citep{mcpspec}. It is part of a broader personal
automation stack that includes local model inference (Ollama on a 24\,GB
Apple-silicon machine), a self-hosted Linux server (Docker, Grafana, Supabase)
and a Tailscale mesh. A single operator cannot justify an enterprise gateway for
this kind of stack, which motivates a lightweight, dependency-free gate.

The HTTP server binds to loopback, and a
\code{cloudflared} tunnel~\citep{cftunnel} is the only path from the Internet
(\cref{fig:arch}). The socket peer is therefore always loopback, so the client
identity is taken from the \code{CF-Connecting-IP} header set at Cloudflare's
edge~\citep{cfheaders}.

\begin{figure}[t]
\centering
\begin{tikzpicture}[
  node distance=3.2mm and 3mm, font=\scriptsize,
  box/.style={draw=accent,rounded corners=2pt,fill=accent!6,align=center,minimum height=6mm,inner sep=2pt},
  gate/.style={box,fill=warn!10,draw=warn},
  arr/.style={-{Stealth[length=1.6mm]},thick,color=black!70}]
\node[box] (net) {Internet clients\\(agents, crawlers,\\scanners, probers)};
\node[box,right=of net] (cf) {Cloudflare edge\\sets \code{CF-Connecting-IP}};
\node[box,right=of cf] (tun) {\code{cloudflared}\\tunnel};
\node[gate,below=5mm of tun] (g) {\sys{} gate\\ban $\to$ bot-UA $\to$ rate};
\node[gate,left=of g] (a) {Bearer auth\\+ auth-fail ban};
\node[gate,left=of a] (f) {Tool filter\\(5 of 29 exposed)};
\node[box,below=4mm of a] (d) {MCP dispatch $\to$ backends};
\draw[arr] (net) -- (cf); \draw[arr] (cf) -- (tun); \draw[arr] (tun) -- (g);
\draw[arr] (g) -- (a); \draw[arr] (a) -- (f); \draw[arr] (f) |- (d);
\begin{scope}[on background layer]
\node[draw=black!40,dashed,rounded corners,fit=(g)(a)(f)(d),inner sep=4pt,
      label={[font=\scriptsize,color=muted]below:loopback-bound origin (\code{127.0.0.1})}] {};
\end{scope}
\end{tikzpicture}
\caption{Deployment. The origin is reachable only through the tunnel, and all
enforcement runs in-process before any tool code executes.}
\label{fig:arch}
\end{figure}

\subsection{Threat model}
Remote, unauthenticated adversaries send arbitrary HTTP requests through
Cloudflare:
(T1) \emph{declared} crawlers and scanners that name themselves in the
\code{User-Agent};
(T2) \emph{undeclared} crawlers that spoof a browser user-agent;
(T3) credential probers using one or many sources;
(T4) flooders;
(T5) path scanners;
(T6) adversaries who rotate addresses, for example within an IPv6 /64.
The goals are to prevent unauthorized tool execution, bound the operator's
resource cost and provide visibility. Out of scope: a compromised operator
machine, Cloudflare itself, volumetric DDoS absorbed at the edge, and attacks
by \emph{authorized} clients such as indirect prompt
injection~\citep{greshake2023indirect}.

\subsection{Guardian v1}
\sys{} v1 (\code{guardian.py}, 263 lines, Python standard library only) makes
one decision per request, before authentication or JSON-RPC parsing:
\begin{enumerate}
\item an unexpired ban returns \code{403};
\item a case-insensitive match against 38 crawler and scanner signatures
      (\code{gptbot}, \code{claudebot}, \code{ccbot}, \code{bytespider},
      \code{sqlmap}, \code{nuclei}, \ldots) bans the client and returns
      \code{403};
\item more than 120 requests in 60\,s returns \code{429}~\citep{rfc6585}
      without a ban.
\end{enumerate}
A bearer token~\citep{rfc6750} is then checked in constant time. Eight
\code{401}s within 300\,s trigger a 60-minute ban, in the style of
\texttt{fail2ban}~\citep{fail2ban} and in line with online-guessing
guidance~\citep{nist80063b}. Bans and a JSONL access log are file-backed so
that a separate stdio process can offer operator tools; rate counters live in
memory. A \textbf{capability filter} removes all 15 \code{local\_only} tools
(offensive recon, breach lookups, the \sys{} controls themselves) and an
operator denylist from \code{tools/list} and \code{tools/call}.

% =====================================================================
\section{Evaluation Method}\label{sec:method}
The harness (\code{artifact/harness2.py}) launches an unmodified server tree
as a separate process on a private port, with a fresh 320-bit bearer token and
a fresh state directory for every scenario. It emulates the Cloudflare hop by
setting \code{CF-Connecting-IP} to documentation-range
addresses~\citep{rfc5737,rfc3849} and replays the scenarios over HTTP/1.1. The
base scenario set (S1--S14) runs against the v1, v2 and v3 trees under the
shipped default policy. Two additional scenarios (S15--S16) exercise
v3-specific features and are skipped on older trees.
Latency is measured end to end on a keep-alive connection: 50 warm-up requests,
then $3\times\num{1000}$ interleaved requests per condition, or 500 per point
for ban-list scaling. Memory is measured with \code{tracemalloc} in a fresh
interpreter. Platform: Apple M5 Pro, macOS 26.6.1, Python 3.12.13. Every number
is read from \code{artifact/results\_v\{1,2,3\}.json}.

% =====================================================================
\section{Guardian v1: Effective, but Keyed on the Wrong Thing}\label{sec:v1}

\finding{What works (RQ1).} On benign traffic, four MCP client user-agents
completed the protocol handshake and 100/100 calls with no false blocks (S1).
A client performing OAuth discovery and an SSE probe before connecting (S1b)
was also unaffected. All 38 declared crawler and scanner signatures were
blocked on the first request, and they stayed blocked after switching to a
benign UA \emph{with the valid token} (S2). A single-source prober was banned
on exactly its 8th failure (S4). A 300-request burst was held to exactly 120
served, with no collateral effect on other clients and recovery after the
window (S6). Over HTTP, \code{tools/list} returned 5 of 29 tools. Each of the
15 \code{local\_only} tools was unlisted and returned JSON-RPC \code{-32601}
when called directly (\code{artifact/surface.json}). The ablation with \sys{}
disabled (S11) served the entire flood and answered every guess.

\finding{What breaks (RQ2).}
\textbf{W1}~Spoofed browser UAs evaded the signature rule 38/38 times (S3);
the token still rejected all of them.
\textbf{W2}~Spreading 700 guesses over 100 sources produced 0 bans and 700
answered guesses (S5).
\textbf{W3}~Each request re-parsed the whole ban file. Median latency rose
from \SI{0.12}{\milli\second} with no bans to \SI{0.54}{}, \SI{3.99}{} and
\SI{21.7}{\milli\second} at \num{1000}, \num{10000} and \num{50000} bans
(\cref{fig:bans}).
\textbf{W4}~Per-IP rate state was never evicted: \SI{87.4}{MB} for
\num{100000} sources ($\sim$\SI{870}{B} each).
\textbf{W5}~Rotating through 200 addresses in one IPv6 /64 created 200 ban
entries (S12), which feeds W3. An end site typically holds at least a
/64~\citep{rfc6177}, so this is cheap for an attacker.
\textbf{W6}~Path scanning (S7) returned only \code{404}/\code{405} and was never
penalized.
\textbf{W7}~Eight mistyped tokens locked a legitimate user out for 60\,min
(S8).
\textbf{W8}~Several smaller gaps (S9--S11). \code{DELETE} bypassed the gate
and was not logged. Banned clients could still make the server read 256\,KB
bodies. On a non-loopback bind, such as the supported tailnet mode, forged
\code{CF-Connecting-IP} headers evaded a ban 5 out of 5 times. With \sys{}
disabled, unenforced ban records were still written.

% =====================================================================
\section{Hardening: Guardian v2}\label{sec:v2}
Each weakness maps to a local change of a few lines (\cref{tab:fixes}); v2
still uses only the standard library.

\begin{table}[t]
\centering\small
\caption{v2 changes, each tied to a measured v1 weakness.}
\label{tab:fixes}
\begin{tabularx}{\columnwidth}{@{}lX@{}}
\toprule
W & Change in v2 \\
\midrule
W3 & Ban list cached in memory; reloaded only when the file's \code{mtime} changes \\
W5 & All state keyed on a \emph{source}: IPv4 address or IPv6 /64 \\
W4 & Idle keys dropped; LRU cap of \num{20000} tracked sources \\
W2 & Global auth-failure alarm (>40 per 300\,s) tightens the per-source threshold from 8 to 3 \\
W7 & Auth-failure bans back off exponentially from 5\,min; optional operator allowlist \\
W6 & Six non-benign \code{404}s per 300\,s ban the source (OAuth discovery and favicon exempt) \\
W8 & Every method gated and logged; gate runs before the body is read; forwarding headers trusted only from a loopback peer; nothing written when disabled \\
\bottomrule
\end{tabularx}
\end{table}

\begin{table*}[t]
\centering\small
\caption{Before/after on identical scenarios (default policy: 120 req/60\,s;
8 auth failures/300\,s; 38 signatures). Bold marks a changed outcome.}
\label{tab:ba}
\begin{tabularx}{\textwidth}{@{}l X X X@{}}
\toprule
\# & Scenario & v1 (deployed) & v2 (hardened) \\
\midrule
S1  & Benign MCP sessions, 4 clients, 100 req & 100/100 served, 0 bans & 100/100 served, 0 bans \\
S1b & Benign client: OAuth discovery + SSE probe, then call & 20$\times$\code{405}; call \code{200}; 0 bans & 16$\times$\code{404}, 4$\times$\code{405}; call \code{200}; 0 bans \\
S2  & 38 declared crawler/scanner UAs & 38/38 blocked on first request; still blocked with valid token & same \\
S3  & Same crawlers, spoofed browser UA & 0/38 caught by UA rule; 38/38 \code{401} & same (not addressable at this layer) \\
S4  & One source, 20 wrong tokens & banned at attempt 8; ban 60\,min & banned at attempt 8; \textbf{ban 5\,min (backs off)} \\
S5  & 100 sources $\times$ 7 wrong tokens & 700 answered, \textbf{0 bans} & \textbf{323 answered, 377 blocked, 95 bans} \\
S6  & 300-request burst, one source & 120 served, 180 \code{429}; recovers at 61\,s & same \\
S7  & Path scan, 10 paths $\times$ \{GET, POST\} & 10$\times$\code{404}, 10$\times$\code{405}; 0 bans & \textbf{banned at 6th \code{404}; 14/20 blocked} \\
S8  & Legit user, 8 typos, then correct token & locked out 60\,min & \textbf{locked out 5\,min} \\
S8b & Same, operator address allowlisted & locked out 60\,min (no allowlist) & \textbf{never banned; call \code{200}} \\
S9  & Banned client: \code{DELETE}; 300\,KB body & \code{200}, unlogged; body read (\code{413}) & \textbf{\code{403}, logged; body unread (\code{403})} \\
S10 & Non-loopback peer forging client IP & \textbf{5/5 forged identities evade ban} & \textbf{0/5; ban keyed on real peer} \\
S12 & 200 bot requests rotating inside one /64 & \textbf{200 ban entries} & \textbf{1 ban}; 20/20 further addresses in /64 blocked; neighbouring /64 unaffected \\
S13 & 4 concurrent clients, 50 req each & 200/200 served, 0 bans & 200/200 served, 0 bans \\
S14 & Long session: 100 req, wait 61\,s, 100 more & 200 served, 0 throttled & 200 served, 0 throttled \\
S11 & Guardian disabled (ablation) & flood served; 1 unenforced ban record & flood served; 0 records \\
\midrule
E2 & Median latency, 0 / 1k / 10k / 50k bans & 0.12 / 0.54 / 3.99 / \textbf{21.7}\,ms & 0.13 / 0.13 / 0.13 / \textbf{0.13}\,ms \\
E2 & Gate overhead (median, empty ban list) & +\SI{5}{\micro\second} (0.119$\to$0.124\,ms) & +\SI{9}{\micro\second} (0.121$\to$0.130\,ms) \\
E3 & Rate-state memory at 1k / 10k / 100k sources & 0.86 / 8.56 / \textbf{87.4}\,MB & 0.91 / 9.01 / \textbf{18.8}\,MB (capped at 20k keys) \\
\bottomrule
\end{tabularx}
\end{table*}

\begin{figure}[t]
\centering
\begin{tikzpicture}
\begin{axis}[
  width=\columnwidth, height=4.4cm,
  xlabel={Active bans}, ylabel={Median latency (ms)},
  xmin=-2000, xmax=52000, ymin=0, ymax=25, scaled x ticks=false,
  xtick={0,10000,20000,30000,40000,50000}, xticklabels={0,10k,20k,30k,40k,50k},
  grid=major, grid style={black!10},
  legend pos=north west, legend style={font=\scriptsize,draw=none,fill=none},
  tick label style={font=\scriptsize}, label style={font=\small}]
\addplot[bad,thick,mark=*,mark size=1.6pt] coordinates
  {(0,0.122) (1000,0.542) (10000,3.986) (50000,21.688)};
\addlegendentry{v1: re-parse per request}
\addplot[good,thick,mark=square*,mark size=1.5pt] coordinates
  {(0,0.129) (1000,0.129) (10000,0.129) (50000,0.128)};
\addlegendentry{v2: mtime-cached}
\end{axis}
\end{tikzpicture}
\caption{Request latency against ban-list size ($n{=}500$ per point). v1
grows linearly (${\approx}\SI{0.43}{\micro\second}$ per entry). v2 is flat,
$169\times$ faster at 50k bans.}
\label{fig:bans}
\end{figure}

\finding{Results (RQ3).} \Cref{tab:ba} gives the full comparison. Every
targeted weakness changed as intended, and the benign scenarios (S1, S1b) were
unchanged: no false blocks, and OAuth-discovery 404s are not treated as
scanning. The costs are small and measurable. v2 adds \SI{4}{\micro\second} of
median overhead over v1. Memory per tracked key rises slightly (${\approx}\SI{940}{B}$
against \SI{870}{B}), but the total is now bounded.

\finding{Trade-offs.} Keying on /64 means that separate customers who share a
/64, which some hosting providers do, share a fate. The global alarm lets a
distributed attacker tighten the threshold for everyone for up to 300\,s. The
worst case is a legitimate user needing 3 attempts instead of 8, which the
allowlist removes. The LRU cap means an attacker with more than \num{20000}
active sources can push out counters. Bans are persisted separately and are
not evicted.

% =====================================================================
\section{Post-Authentication Extension: Guardian v3}\label{sec:v3}
v1 and v2 treat the bearer token as monolithic: one token grants access to
every public tool. Once the token leaks---from a client's configuration, an
agent's context window, or a log line---the holder has the full tool surface.
And once authenticated, a client faces only the global rate limit, so a
misbehaving or compromised agent can call one tool hundreds of times without
triggering throttling.

\subsection{Scoped tokens}
v3 replaces the single token with a set of named tokens, each scoped to a list
of tool names (or \code{["*"]} for full access). The legacy single token remains
accepted as a full-access credential for backward compatibility.
\code{tools/list} returns only the tools the presenting token can reach, and
\code{tools/call} to an out-of-scope tool returns JSON-RPC \code{-32601}
without executing.

\subsection{Per-tool rate limiting}
v3 adds a per-(source, tool) sliding window on top of the global rate limit.
The default is 30 calls per 60\,s per tool per source. Exceeding it returns
JSON-RPC \code{-32005}; other tools and other sources are unaffected.
Allowlisted sources are exempt, consistent with the admission gate.

\finding{Results.}
S15 verifies token scoping: a full-access token sees all 5 public tools; a
scoped token sees only its 2 allowed tools (\code{devops\_status},
\code{devops\_health}); calling \code{devops\_logs} with the scoped token
returns \code{-32601}. S16 verifies per-tool rate limiting: with a limit of
10 calls per 60\,s, the first 10 calls to \code{shell\_explain} are served
and the next 20 return \code{-32005}; a different tool
(\code{shell\_suggest}) from the same source is unaffected.

\finding{Cost.} Both features add no measurable per-request latency:
they are checked inside \code{\_handle\_filtered} after authentication
succeeds, using the same sliding-window and LRU data structures as the
admission gate. Token lookup is a linear scan of the token set (typically
2--5 entries) with constant-time comparison per token. Total added code:
32 lines across three files.

% =====================================================================
\section{Real-World Exposure}\label{sec:real}
\noindent\textbf{Method.} We exposed the endpoint running v2 at its public
hostname for a fixed window. The server logged, for every request, the
timestamp, client address, user-agent, method, path, status, the outcome of
the authentication check and the declared body length. These are enough to
re-decide each request under any policy. The analysis tool
(\code{artifact/analyze\_trace.py}) replays the trace in timestamp order on
the trace's own clock through three policies: no gate (auth only), v1 and v2.
Each runs its own unmodified code tree with isolated state. We validated the
replay against ground truth: on 1211 requests from the v2 harness logs, it
reproduced the live server's status code for every request (1211/1211,
\code{artifact/replay\_fidelity.json}). Operator traffic was excluded by
user-agent. Only aggregates are published.

\ifrealdata
\RTCompositionTable
\RTPolicyTable
\noindent\textbf{Results (RQ4).} Over \RTHours\,h the endpoint received
\RTRequests{} requests from \RTSources{} sources (\RTvSix{} IPv6 /64s). The
ten busiest sources sent \RTTopTen{} of the traffic (\cref{tab:realcomp}).
Replayed through each policy (\cref{tab:realpolicy}), v1 blocked
\RTvOneBlocked{} requests at the gate and v2 blocked \RTvTwoBlocked{}. The
share of non-MCP-client requests stopped before authentication was
\RTvOneUnwantedPct{} and \RTvTwoUnwantedPct{} respectively. With no gate,
\RTNoneAuth{} requests reached the authentication layer, compared with
\RTvOneAuth{} under v1 and \RTvTwoAuth{} under v2.
The composition and policy comparison are discussed in \cref{tab:realcomp,tab:realpolicy}.
\else
The exposure study is in progress at the time of writing; results will be
incorporated when the observation window closes. The methodology and replay
tooling are complete and validated (1211/1211 exact status match on harness
traffic).
\fi

% =====================================================================
\section{Related Work}

\finding{MCP threat landscape.}
Hou et al.~\citep{hou2025mcp} give the first survey of MCP security threats and
future directions. Javed et al.~\citep{mcpthreatmodel2026} apply STRIDE and
DREAD threat modeling across five MCP components and identify tool poisoning as
the most prevalent client-side vulnerability. The NSA's AI Security Center
published the first US government guidance on MCP
security~\citep{nsamcp2026}, warning that adoption has outpaced safeguards and
recommending controls including input screening, least-privilege interfaces,
audit logging, and overload protection. The Cloud Security Alliance provides
complementary practitioner guidance~\citep{csamcp2026}. OWASP catalogues
broader LLM application risks~\citep{owaspllm2025}, and Bajpai
et al.~\citep{mcpsafetyaudit2025} demonstrate major exploits against MCP-enabled
LLM systems.

\finding{Post-authentication attacks.}
The largest body of MCP security research studies what an authorized or
compromised client can do once inside. MCPTox~\citep{mcptox2025} benchmarks
tool poisoning across 45 live servers and 353 tools, finding attack success
rates above 72\%. MCP-ITP~\citep{mcpitp2026} automates implicit tool poisoning
generation in a black-box setting. TrustShiftProbe~\citep{trustshift2026}
characterizes staged trust-escalation attacks. Indirect prompt
injection~\citep{greshake2023indirect} targets authorized data paths. Our
threat model explicitly excludes all of these: we study the \emph{pre-authentication}
admission layer, not the post-authentication trust boundary.

\finding{Defense mechanisms.}
CASCADE~\citep{cascade2026} proposes a three-tiered local defense for
prompt-injection detection in MCP systems (regex and entropy pre-filtering,
semantic analysis, output filtering), evaluated on 5{,}000 samples at 95.85\%
precision. Tamayo~\citep{tamayo2026slr} surveys 73 studies and categorizes
existing MCP defenses into static analysis, runtime guardrails, middleware
access control and formal verification.
None of the 73 surveyed studies measures an admission-layer defense for an
Internet-exposed MCP endpoint.

\finding{Established network defense.}
Log-driven banning~\citep{fail2ban} and throttling of online
guessing~\citep{nist80063b} are long established. Crawler consent rests on
convention~\citep{rfc9309}, and restrictions aimed at AI crawlers have grown
sharply~\citep{longpre2024consent}. Our S3 result confirms that
convention-based identification cannot be relied on.

\finding{Gap.}
Prior work on MCP security either catalogues threats, studies post-authentication
attacks, or prescribes controls without empirical measurement. We contribute a
measured, before/after evaluation of admission-layer defense on a deployed,
Internet-facing MCP endpoint, show where per-address state fails, and quantify
the cost of fixes on identical workloads.

% =====================================================================
\section{Limitations}
This is a single-system, single-operator study. Controlled experiments ran on
loopback with an emulated Cloudflare hop, so edge behaviour is assumed rather
than measured, and the latency figures exclude network and tunnel time. The
benign workload now covers concurrent multi-client sessions and rate-window
recovery (S13, S14) but not SSE streaming or mixed read/write tool workloads.
The v2 defaults (the /64 prefix, the 20k cap and the global-alarm thresholds)
and v3 defaults (30 calls per tool per window) are reasonable choices, not
tuned optima.
The real-traffic window covers a single hostname whose discoverability depends
on certificate-transparency logs and DNS enumeration by third parties. The
techniques evaluated (rate limiting, UA signature blocking, auth-failure banning,
token scoping)
are individually well established~\citep{fail2ban,nist80063b}; the contribution
is their measured application and failure modes in the MCP context, not the
techniques themselves.

\section{Conclusion}
A few hundred lines of dependency-free code in front of an MCP endpoint,
together with a strong token and a minimized tool surface, measurably stop
declared crawlers, credential probing, path scanning and floods at
microsecond cost. The first version failed where its state was keyed on
individual addresses and allowed to grow without bound. Distributed and
rotating adversaries slipped through, and a growing ban list turned the
defense into a latency liability. Keying on network sources and bounding
state fixed each measured failure without new false blocks. Extending
defense past authentication with scoped tokens and per-tool rate limits
bounds the blast radius of a leaked credential and prevents authorized
abuse of individual tools. Self-hosted agent endpoints should treat such
gates as cost and visibility controls, and never as the authentication
boundary.

\smallskip\noindent\textbf{Artifact.} \code{artifact/}: \code{harness2.py},
\code{harness.py} and \code{surface.py} (tool surface), \code{analyze\_trace.py}, \code{gen\_realtraffic\_tex.py},
\code{results\_v1.json}, \code{results\_v2.json}, \code{results\_v3.json}, \code{surface.json},
\code{replay\_fidelity.json}, and aggregate \code{real\_traffic.json};
\url{https://github.com/TGvenomYT/agent-endpoint-defense}.

\smallskip\noindent\textbf{Ethics.} Controlled experiments targeted only the
author's own software, using documentation-range addresses. The public
exposure passively logged requests sent to the author's own hostname. No
third-party system was probed and no traffic source was attacked. Raw logs,
which contain client IP addresses, are kept private and deleted after
analysis. Only aggregates are released.

\balance
\small
\bibliographystyle{plainnat}
\bibliography{refs}

\end{document}
